\section{Balancing Diversity with Coherence: Top-p Filtering}
\label{sec:top-p-filtering}

% WHAT TO WRITE -------------------------------------------------------
% A status snapshot: what is done, in progress, and planned.
% ---------------------------------------------------------------------

High temperature settings can occasionally cause the model to sample "weird" or nonsensical tokens (e.g., sampling "mistress" when the query is "The capital of Germany is..."). To prevent this, Top-p Filtering (also known as Nucleus Sampling) is implemented.

\subsection{The Four-Step Top-p Process}

This filter ensures that the model only samples from the most plausible tokens:

\begin{enumerate}
    \item \textbf{Sort:} Arrange token probabilities in descending order.
    \item \textbf{Cumulative Sum:} Calculate the running total of probabilities.
    \item \textbf{Apply Threshold:} Select the smallest subset of tokens whose cumulative probability exceeds the threshold $p (e.g., p=0.9)$.
    \item \textbf{Renormalize:} Rescale the remaining probabilities so they sum to 1.0, creating a new, valid distribution for sampling.
\end{enumerate}

By dropping low-probability tokens, the model maintains coherence while still allowing for the variability needed for advanced inference scaling strategies.